\section*{第 11 章：简单扩展}
\addcontentsline{toc}{section}{第 11 章：简单扩展}

\begin{taplsolutionentrybox}
\phantomsection
\label{sol:translator:11.2-unit-puzzle}
\textbf{第 11.2 节未编号练习}

答案是可数无限多个。事实上，即使把类型固定为
\[
(\tapltype{Unit}\to\tapltype{Unit})
\to\tapltype{Unit}\to\tapltype{Unit}
\]
也能构造无限多个不同的闭范式。对每个 \(n\in\Nat\)，定义
\[
c_n
=\lambda f:\tapltype{Unit}\to\tapltype{Unit}.\,
 \lambda x:\tapltype{Unit}.\,f^n x
\]
其中 \(f^0x=x\)，而 \(f^{n+1}x=f(f^n x)\)。省略上面已经给出的类型标注后，
开头几项是
\[
\begin{aligned}
c_0&=\lambda f.\,\lambda x.\,x,\\
c_1&=\lambda f.\,\lambda x.\,f\,x,\\
c_2&=\lambda f.\,\lambda x.\,f(f\,x),\\
c_3&=\lambda f.\,\lambda x.\,f(f(f\,x))
\end{aligned}
\]
这些项只有变量 \(f\) 和 \(x\)，且二者都已被 lambda 抽象绑定，所以它们
都是闭项。每个应用的函数位置都是变量 \(f\)，而不是 lambda 抽象，因此没有
\(\beta\) 可归约式，所有 \(c_n\) 都是范式。不同的 \(n\) 对应不同层数的
\(f\) 应用，所以这些范式在语法上两两不同。这说明闭范式至少有可数无限多个；
另一方面，每个 lambda 项都是有限字符串，所有项的集合至多可数，因此闭范式
恰好有可数无限多个。

这组项正是以 \(\tapltype{Unit}\) 为底层类型写出的 Church 数。尽管它们在
语法上不同，但在没有副作用的 \(\tapltype{Unit}\) 语言中，给定函数和参数后，
最终可观察到的结果仍只能是 \(\mathsf{unit}\)。因此，这个例子也展示了
“语法不同”与“可观察行为不同”并不是一回事。

\taplsolutionbacklink{ex:11.2-unit-puzzle}{返回第 11.2 节脚注谜题}
\end{taplsolutionentrybox}

\begin{taplsolutionentrybox}
\phantomsection
\label{sol:translator:11.8.1}
\textbf{练习 11.8.1}

可把记录投影规则显式写成
\[
\inferrule*[right=\rulename{E-ProjRcd}]
  {j\in1..n \\
   l=l_j}
  {\{l_i=v_i\}^{i\in1..n}.l\trans v_j}
\]
这里还隐含记录标签两两不同这一良构条件。与正文中的简写相比，这一版本把
“所投影标签 \(l\) 恰好是第 \(j\) 个字段的标签”写进了前提，因此无需再靠
文字约定解释下标 \(j\) 从何而来。

\taplsolutionbacklink{ex:11.8.1}{返回练习 11.8.1}
\end{taplsolutionentrybox}
